
\documentclass{article}
\usepackage{colortbl}
\usepackage{makecell}
\usepackage{multirow}
\usepackage{supertabular}

\begin{document}

\newcounter{utterance}

\centering \large Interaction Transcript for game `eqbench', experiment `eq\_bench', episode 30 with reflexion{-}llama{-}3.1{-}8b{-}instruct.
\vspace{24pt}

{ \footnotesize  \setcounter{utterance}{1}
\setlength{\tabcolsep}{0pt}
\begin{supertabular}{c@{$\;$}|p{.15\linewidth}@{}p{.15\linewidth}p{.15\linewidth}p{.15\linewidth}p{.15\linewidth}p{.15\linewidth}}
    \# & \multicolumn{2}{c}{Player} && \multicolumn{2}{c}{Game Master} \\
    \hline

    \theutterance \stepcounter{utterance}  
    & & & \multicolumn{4}{p{0.6\linewidth}}{
        \cellcolor[rgb]{0.9,0.9,0.9}{
            \makecell[{{p{\linewidth}}}]{
                \texttt{\tiny{[P1$\langle$GM]}}
                \texttt{You have attempted this same task before and did not succeed. These are the plans you wrote after those attempts.} \\
\texttt{Trial \#1: New plan: } \\
\\ 
\texttt{Upon re{-}examining the dialogue, I realized that I focused too much on Jamal's responses and didn't consider Sophia's perspective. I should have taken into account her concerns about the project's timeline and budget, as well as her frustration with Jamal's unconventional approach. } \\
\\ 
\texttt{To improve my analysis, I will pay closer attention to Sophia's emotional cues and consider how Jamal's responses might be affecting her. I will also look for moments where Jamal's words and actions might be perceived as dismissive or uncooperative, which could contribute to Sophia's feelings of resentment. } \\
\\ 
\texttt{Additionally, I will consider the context of the conversation and the power dynamics at play. As the project lead, Sophia has a vested interest in the project's success, and Jamal's defiance might be perceived as a threat to her authority. } \\
\\ 
\texttt{By taking a more nuanced approach and considering multiple perspectives, I hope to provide a more accurate analysis of Jamal's emotional state.} \\
\\ 
\\ 
\\ 
\texttt{Your task is to predict the likely emotional responses of a character in this dialogue:} \\
\\ 
\texttt{Sophia: Jamal, I've been looking at your design drafts. They're... unconventional.} \\
\texttt{Jamal: That's one way to put it, Soph. I prefer to call it 'innovative'.} \\
\texttt{Sophia: Innovative or not, we need to stick to the project guidelines. We can't afford to take risks.} \\
\texttt{Jamal: Risks? That's what creativity is all about, Sophia. We need to break the mold.} \\
\texttt{Sophia: This isn't about breaking molds, Jamal. This is about delivering a project on time and within budget.} \\
\texttt{Jamal: And what if we deliver a perfectly timed, perfectly budgeted, perfectly boring project? What then?} \\
\texttt{Sophia: We're not here to entertain, Jamal. We're here to work.} \\
\texttt{Jamal: And who said work can't be entertaining? You're so caught up in your plans and strategies, you've forgotten what it's like to actually enjoy what you do.} \\
\texttt{Sophia: This isn't about enjoyment. This is about responsibility.} \\
\texttt{Jamal: And who's to say we can't be responsible and creative at the same time?} \\
\texttt{Sophia: I'm not saying that, Jamal. I'm saying we need to prioritize.} \\
\texttt{Jamal: And I'm saying we need to think outside the box.} \\
\texttt{Sophia: We don't have time for this. I'm going to revise the project plan.} \\
\texttt{Jamal: And I'm going to continue working on my designs.} \\
\texttt{Sophia: Fine. But remember, Jamal, if this project fails, it's on you.} \\
\texttt{Jamal: And if it succeeds, it's on me too, right?} \\
\texttt{Sophia: We'll see about that.} \\
\texttt{{[}End dialogue{]}} \\
\\ 
\texttt{At the end of this dialogue, Jamal would feel...} \\
\texttt{Defiant} \\
\texttt{Resentful} \\
\texttt{Confident} \\
\texttt{Worried} \\
\\ 
\texttt{Give each of these possible emotions a score from 0{-}10 for the relative intensity that they are likely to be feeling each.} \\
\\ 
\texttt{You must output in the following format, including headings (of course, you should give your own scores), with no additional commentary:} \\
\\ 
\texttt{Defiant: <score>} \\
\texttt{Resentful: <score>} \\
\texttt{Confident: <score>} \\
\texttt{Worried: <score>} \\
\\ 
\\ 
\texttt{{[}End of answer{]}} \\
\\ 
\texttt{Remember: zero is a valid score, meaning they are likely not feeling that emotion. You must score at least one emotion > 0.} \\
\\ 
\texttt{Your answer:} \\
            }
        }
    }
     \\ \\

    \theutterance \stepcounter{utterance}  
    & \multicolumn{4}{p{0.6\linewidth}}{
        \cellcolor[rgb]{0.561,0.737,0.561}{
            \makecell[{{p{\linewidth}}}]{
                \texttt{\tiny{[P1$\rangle$GM]}}
                \texttt{Defiant: 8} \\
\texttt{Resentful: 6} \\
\texttt{Confident: 4} \\
\texttt{Worried: 0} \\
            }
        }
    }
    & & \\ \\

    \theutterance \stepcounter{utterance}  
    & & & \multicolumn{2}{p{0.3\linewidth}}{
        \cellcolor[rgb]{0.9,0.9,0.9}{
            \makecell[{{p{\linewidth}}}]{
                \texttt{\tiny{[GM$|$GM]}}
                \texttt{Defiant: 8} \\
\texttt{Resentful: 6} \\
\texttt{Confident: 4} \\
\texttt{Worried: 0} \\
            }
        }
    }
    & & \\ \\

    \theutterance \stepcounter{utterance}  
    & & & \multicolumn{2}{p{0.3\linewidth}}{
        \cellcolor[rgb]{0.9,0.9,0.9}{
            \makecell[{{p{\linewidth}}}]{
                \texttt{\tiny{[GM$|$GM]}}
                \texttt{Defiant: 8} \\
\texttt{Resentful: 3} \\
\texttt{Confident: 7} \\
\texttt{Worried: 0} \\
            }
        }
    }
    & & \\ \\

    \theutterance \stepcounter{utterance}  
    & & & \multicolumn{2}{p{0.3\linewidth}}{
        \cellcolor[rgb]{0.9,0.9,0.9}{
            \makecell[{{p{\linewidth}}}]{
                \texttt{\tiny{[GM$|$GM]}}
                \texttt{game\_result = LOSE} \\
            }
        }
    }
    & & \\ \\

    \theutterance \stepcounter{utterance}  
    & & \multicolumn{4}{p{0.6\linewidth}}{
        \cellcolor[rgb]{0.95,0.95,0.95}{
            \makecell[{{p{\linewidth}}}]{
                \texttt{\tiny{[P1$\leftrightarrow$P]}}
                \texttt{Post{-}episode Reflexion (not sent to GM):} \\
\texttt{New plan: } \\
\\ 
\texttt{Upon re{-}examining the dialogue, I realized that I focused too much on Jamal's responses and didn't consider Sophia's perspective. I should have taken into account her concerns about the project's timeline and budget, as well as her frustration with Jamal's unconventional approach. } \\
\\ 
\texttt{To improve my analysis, I will pay closer attention to Sophia's emotional cues and consider how Jamal's responses might be affecting her. I will also look for moments where Jamal's words and actions might be perceived as dismissive or uncooperative, which could contribute to Sophia's feelings of resentment. } \\
\\ 
\texttt{Additionally, I will consider the context of the conversation and the power dynamics at play. As the project lead, Sophia has a vested interest in the project's success, and Jamal's defiance might be perceived as a threat to her authority. } \\
\\ 
\texttt{To accurately predict Jamal's emotional state, I will ask myself the following questions:} \\
\\ 
\texttt{1. How does Jamal's defiance affect Sophia's emotional state?} \\
\texttt{2. Are there moments where Jamal's words and actions might be perceived as dismissive or uncooperative?} \\
\texttt{3. How does the power dynamic between Jamal and Sophia influence the conversation?} \\
\texttt{4. Are there any moments where Jamal shows vulnerability or uncertainty?} \\
\\ 
\texttt{By taking a more nuanced approach and considering multiple perspectives, I hope to provide a more accurate analysis of Jamal's emotional state.} \\
            }
        }
    }
    & \\ \\

\end{supertabular}
}

\end{document}
